\section{Signals and probes}\label{sec:signals}

All internal signals are read from the forward pass that produced the proposal. The stored answer is re-run once
with teacher forcing to extract attention and hidden states, and the agreement between the teacher-forced
argmax and the generated token is checked. Agreement was 96--99.9\% for the bfloat16 models, except
Qwen2.5-1.5B in path planning (92.8\%). It was lower for the 4-bit Qwen2.5-7B (84\% in path planning, 91\% in the
CSTR), because quantised kernels make the two passes diverge more often. Signals are taken at four
relative depths $\mathcal L=\{0.25,0.5,0.75,1.0\}$ of the network, where depth $1.0$ is the last layer. Let
$\mathcal T$ denote the answer tokens and $\mathcal P$ the prompt tokens. Let $\alpha^{(\ell,\eta)}_{j,i}$ be the
attention weight of head $\eta$ in layer $\ell$ from the position that predicts answer token $j$ to prompt token
$i$, and let $\bar\alpha^{(\ell)}_{j,i}$ be its mean over heads.

\subsection{Internal signals}

\paragraph{Token confidence} For answer token $j$ with predictive distribution $p_j$ and selected token $y_j$,
we compute the log-probability $\log p_j(y_j)$, the entropy $-\sum_v p_j(v)\log p_j(v)$ and the margin between
the two most probable tokens. The features are their mean, minimum and maximum over all answer tokens and over
the numeric tokens. This is the cheapest internal signal and needs no extra pass.

\paragraph{Attention by prompt region} The prompt is partitioned into coarse regions $R$: system instructions,
plant knowledge or graph description, task state, and user instructions. The mass that answer token $j$ places on
region $R$ in layer $\ell$ is
\begin{equation}\label{eq:regionmass}
A^{(\ell)}_j(R)=\frac{\sum_{i\in R}\bar\alpha^{(\ell)}_{j,i}}{\sum_{i\in\mathcal P}\bar\alpha^{(\ell)}_{j,i}},
\end{equation}
averaged over answer tokens. Head-maximum variants and the attention entropy are included as well. The signal
asks \emph{where in the prompt} the model looked while answering.

\paragraph{Hidden states} The residual-stream vectors at the last prompt token and their mean over answer tokens,
at each depth in $\mathcal L$, are standardised. They are reduced to 64 principal components fitted on the
training partition only. Linear probes on such representations are known to recover correctness information
\citep{alain2016probes,azaria2023internal,orgad2025llms}.

\paragraph{Region-based attention grounding} Each answer contains \emph{decision units}: a step $u\to v$ of a
path, or one of the three setpoints. For unit $u$ with decision tokens $\mathcal T_u\subseteq\mathcal T$, a
\emph{deciding region} $R^{\star}_u$ is fixed in advance. It is the part of the prompt that contains the
information on which the unit's correctness depends. The grounding share of unit $u$ is
\begin{equation}\label{eq:grounding}
g_u=\frac{1}{|\mathcal L|}\sum_{\ell\in\mathcal L}\;\frac{1}{|\mathcal T_u|}\sum_{j\in\mathcal T_u}
\frac{\sum_{i\in R^{\star}_u}\bar\alpha^{(\ell)}_{j,i}}{\sum_{i\in R_{\mathrm{task}}}\bar\alpha^{(\ell)}_{j,i}},
\end{equation}
where $R_{\mathrm{task}}$ contains all task-content tokens (all adjacency lines in path planning; all system,
knowledge-graph and snapshot tokens in the CSTR). Template tokens and previously generated tokens are excluded,
so $g_u$ is the fraction of the model's attention to task content that falls on the content that matters. The
answer-level features are the minimum $g_{\min}=\min_u g_u$, which was pre-registered as primary, the mean over
units, per-layer variants, and, in path planning, the fraction of steps for which $R^{\star}_u$ is the
most-attended line.

The deciding regions encode domain knowledge and were fixed before any grounding value was computed. In path
planning a step $u\to v$ is legal if and only if $v$ appears in the adjacency line of $u$, so $R^{\star}_u$ is that
line. The definition is exact. In the CSTR the deciding regions of each setpoint were taken from the process
structure stated in the prompt (Table~\ref{tab:regions}). They are the measurements that indicate the setpoint's
effect, the current setpoints, the knowledge-graph statement of the loop that the setpoint drives, and the
statement of the mass and energy balances (operator PO5). Figure~\ref{fig:grounding} illustrates the intended
reading of the signal and the two failure types it is meant to separate.

\begin{table}[!htbp]
\centering
\small
\caption{Pre-registered deciding regions $R^{\star}$ per CSTR setpoint (prompt v3.1). Snapshot fields are single
lines of the prompt's plant snapshot; knowledge-graph (KG) statements are complete Turtle statements that mention
the operator.}
\label{tab:regions}
\resizebox{\textwidth}{!}{%
\begin{tabular}{ll}
\toprule
Setpoint & Deciding regions \\
\midrule
Feed flow $F_{\mathrm{in,sp}}$ & $u_{\mathrm{cool}}$, $T_{\mathrm{meas}}$, current setpoints, KG PO8 (flow loop), KG PO5 (balances) \\
Level $L_{\mathrm{sp}}$ & $L_{\mathrm{meas}}$, $u_{\mathrm{pump}}$, current setpoints, KG PO6 (level loop), KG PO5 \\
Temperature $T_{\mathrm{sp}}$ & $T_{\mathrm{meas}}$, $u_{\mathrm{cool}}$, current setpoints, KG PO7 (temperature loop), KG PO5 \\
\bottomrule
\end{tabular}}
\end{table}

\begin{figure}[!htbp]
\centering
% Region-based attention grounding: schematic of a retrieval failure (path planning) and a reasoning failure (CSTR).
% Bar lengths are illustrative, not measured values.
\resizebox{\textwidth}{!}{%
\begin{tikzpicture}[font=\footnotesize,
  row/.style={anchor=west, font=\ttfamily\scriptsize, inner sep=1pt},
  hdr/.style={font=\scriptsize\bfseries, align=center},
]
\def\bw{2.4}  % bar length for an attention share of 1
% ------------------------------------------------------------------ (a) path planning
\node[anchor=west, font=\small\bfseries] at (-0.2,0.85) {(a) Path planning: a retrieval failure};
\node[anchor=west, font=\scriptsize, align=left] at (-0.2,0.3) {query $2\to 6$; writing the step after node 0};
\node[hdr] at (0.7,-0.25) {adjacency\\[-2pt]lines};
\node[hdr] at (2.75,-0.25) {legal step $0\!\to\!3$};
\node[hdr] at (5.35,-0.25) {illegal step $0\!\to\!5$};
\foreach \i/\txt/\a/\b in {1/{0: 3, 7}/0.52/0.07, 2/{1: 4, 6}/0.05/0.06, 3/{2: 0, 5}/0.10/0.43, 4/{3: 1, 6}/0.07/0.08,
                          5/{4: 2}/0.04/0.05, 6/{5: 6}/0.06/0.14, 7/{6: 4}/0.08/0.09, 8/{7: 1}/0.08/0.08} {
  \pgfmathsetmacro{\y}{-0.45*\i-0.5}
  \ifnum\i=1 \fill[cGrd!22] (-0.15,\y-0.2) rectangle (6.95,\y+0.2); \fi
  \node[row] at (0,\y) {\txt};
  \fill[cInt!80] (1.75,\y-0.12) rectangle ++(\a*\bw,0.24);
  \fill[cInt!80] (4.35,\y-0.12) rectangle ++(\b*\bw,0.24);
}
\draw[black!40] (1.75,-0.65) -- (1.75,-4.3);
\draw[black!40] (4.35,-0.65) -- (4.35,-4.3);
\node[anchor=west, font=\scriptsize, text=cGrd!70!black] at (-0.15,-4.55) {deciding region $R^{\star}$: the source node's line};
\node[font=\scriptsize] at (2.75,-5.0) {$g=0.52$ (high)};
\node[font=\scriptsize] at (5.35,-5.0) {$g=0.07$ (low)};
\node[font=\scriptsize, align=center, text=black!70] at (3.4,-5.65) {low grounding flags the failure:\\[-1pt]the model did not consult the line that decides the step};

% ------------------------------------------------------------------ (b) CSTR
\begin{scope}[xshift=8.0cm]
\node[anchor=west, font=\small\bfseries] at (-0.2,0.85) {(b) CSTR: a reasoning failure};
\node[anchor=west, font=\scriptsize, align=left] at (-0.2,0.3) {reactor hot, cooling saturated; writing $F_{\mathrm{in,sp}}$};
\node[hdr] at (1.0,-0.25) {prompt\\[-2pt]regions};
\node[hdr] at (3.45,-0.25) {passing proposal\\[-2pt](cuts feed)};
\node[hdr] at (6.05,-0.25) {failing proposal\\[-2pt](raises feed)};
\foreach \i/\txt/\d/\a/\b in {1/{T\_meas}/1/0.13/0.14, 2/{u\_cool}/1/0.12/0.13, 3/{setpoints}/1/0.11/0.10,
                             4/{KG: PO8 (flow)}/1/0.07/0.07, 5/{KG: PO5 (balances)}/1/0.05/0.05,
                             6/{L\_meas}/0/0.05/0.04, 7/{u\_pump}/0/0.04/0.04, 8/{KG: PO7 (temp.)}/0/0.08/0.08} {
  \pgfmathsetmacro{\y}{-0.45*\i-0.5}
  \ifnum\d=1 \fill[cGrd!22] (-0.15,\y-0.2) rectangle (7.65,\y+0.2); \fi
  \node[row] at (0,\y) {\txt};
  \fill[cInt!80] (2.75,\y-0.12) rectangle ++(\a*\bw*2,0.24);
  \fill[cInt!80] (5.05,\y-0.12) rectangle ++(\b*\bw*2,0.24);
}
\draw[black!40] (2.75,-0.65) -- (2.75,-4.3);
\draw[black!40] (5.05,-0.65) -- (5.05,-4.3);
\node[anchor=west, font=\scriptsize, text=cGrd!70!black] at (-0.15,-4.55) {deciding regions $R^{\star}$ of $F_{\mathrm{in,sp}}$ (pre-registered)};
\node[font=\scriptsize] at (3.45,-5.0) {$g=0.48$};
\node[font=\scriptsize] at (6.05,-5.0) {$g=0.49$};
\node[font=\scriptsize, align=center, text=black!70] at (3.9,-5.65) {grounding cannot flag the failure:\\[-1pt]the model reads the right regions and draws the wrong conclusion};
\end{scope}
\end{tikzpicture}}

\caption{Region-based attention grounding (schematic; bar lengths are illustrative). Each bar is the share of
the model's attention to task content that falls on a prompt region while the model writes one decision.
(a)~In path planning the step after node~0 is legal only if its target appears in line ``0: 3, 7''. A legal step
attends to that line (high $g$). An illegal step attends elsewhere, here to line ``2: 0, 5'', which contains the
wrongly chosen node (low $g$). Low grounding flags the \emph{retrieval} failure. (b)~In the CSTR, a proposal that
cuts the feed and one that raises it while cooling is saturated both attend to the deciding regions. The failing
proposal is a \emph{reasoning} failure, and grounding cannot distinguish it.}
\label{fig:grounding}
\end{figure}

\subsection{Observable signals (baselines)}
The engineering baseline uses only information that is visible without opening the model.
\begin{itemize}
\item \textbf{Path planning:} the numbers of nodes and edges, the edge density, the length of the shortest path
(computed by breadth-first search on the prompt's graph), the length of the proposed path, and indicators of
whether the path has two nodes, repeats a node or has the correct endpoints. Two of these features use part of
the validator's logic (a shortest-path computation and an endpoint check). This makes the observable baseline
deliberately strong.
\item \textbf{CSTR:} the measured temperature and level, the three actuator commands, the anomaly ratio, the
control zone, the alarm flags, and the proposed setpoint changes (absolute and, for the feed, relative). A
\emph{plant-readings-only} variant omits the action features. Hidden plant parameters, the fault family, its
severity and onset time are never features.
\end{itemize}
The observable CSTR screen is in effect a learned surrogate of the digital twin: it learns how the plant responds
to a change from a state. The internals-only screen is not. It sees no plant reading and no proposed setpoint,
only the agent's forward pass.

\subsection{Signal sets and probes}
We compare ten signal sets in path planning and eleven in the CSTR, each fitted as a single probe: observables;
plant readings only (CSTR); token confidence; attention by region; region grounding; hidden states; \emph{all
internals} (token confidence, attention and hidden states); all internals with grounding, called
\emph{internals only} below; observables with grounding; observables with all internals; and observables with all
internals and grounding, called \emph{combined}. Every probe is an $\ell_2$-regularised logistic regression fitted
on the training partition, with median imputation and standardisation fitted on the same partition, and
Platt-calibrated \citep{platt1999} on a separate calibration partition. Logistic probes were chosen deliberately.
They are cheap, they are standard for probing \citep{alain2016probes,belinkov2022probing}, and a fixed model
class keeps the comparison between signal sets fair. Any advantage of the internal signals must come from the
signals, not from a more flexible probe.
